\chapter*{چکیده}
\addcontentsline{toc}{chapter}{چکیده}

در این پروژه مقالهٔ «یک شبکهٔ سادهٔ ترنسفورمرگونه برای ابرتفکیک‌پذیری سبک تصویر» را از ابتدا بررسی و بازپیاده‌سازی کردیم. مسئله این است که از یک تصویر کم‌وضوح، تصویری با ابعاد بزرگ‌تر و جزئیات نزدیک به نمونهٔ مرجع ساخته شود؛ در عین حال مدل باید برای کاربردهای کم‌تاخیر سبک بماند. ایدهٔ اصلی مقاله جایگزینی خودتوجهی پرهزینه با «پیمانه‌سازی کانولوشنی» است. در این عملگر، دو شاخهٔ کانولوشنی ساخته می‌شود، خروجی آن‌ها عضو‌به‌عضو ضرب می‌شود و در پایان یک کانولوشن $3\times3$ ویژگی محلی را استخراج می‌کند. این عملگر در ساختاری شبیه بلوک ترنسفورمر، همراه با نرمال‌سازی لایه‌ای، پرسپترون چندلایه و اتصال مانده قرار می‌گیرد.

پیاده‌سازی ما شامل آماده‌سازی زوج‌های کم‌وضوح و پروضوح، افزایش داده، مدل کامل، آموزش با زیان $L_1$، ارزیابی روی کانال روشنایی، ذخیرهٔ نقاط وارسی، شروع گرم، رسم نمودار و آزمایش‌های حذف مؤلفه است. جزئیاتی که در مقاله صریح نبودند با تطبیق شمار پارامترها استنتاج شدند؛ مدل مقیاس دو دقیقاً $881{,}912$ پارامتر دارد و شمار پارامتر همهٔ نسخه‌های حذف مؤلفه نیز با جدول مقاله برابر شد. هر سه مدل روی همهٔ ۸۰۰ تصویر آموزشی DIV2K اجرا شدند. برای هر مقیاس، ۱۰۰ دوره آموزش از صفر و سپس ۱۰۰ دوره آموزش با شروع گرم انجام دادیم و خروجی پایان مرحلهٔ دوم را روی همهٔ تصاویر پنج مجموعهٔ معیار سنجیدیم. در مقیاس دو، نتیجهٔ Set5 برابر ۳۷٫۱۶۳ دسی‌بل و ۰٫۹۵۷۵ و نتیجهٔ Urban100 برابر ۳۰٫۳۴۳ دسی‌بل و ۰٫۹۰۹۳ شد. شروع گرم در همهٔ ۱۵ حالت، PSNR را نسبت به مرحلهٔ نخست افزایش داد و فاصلهٔ مدل با درون‌یابی دو‌مکعبی نیز در همهٔ مجموعه‌ها مثبت بود. باقی‌ماندن فاصله با اعداد مقاله با بودجهٔ ۲۰۰‌دوره‌ای ما در برابر آموزش ۱۰۰۰‌دوره‌ای و بازآموزی مقاله سازگار است.

\textbf{واژه‌های کلیدی:} ابرتفکیک‌پذیری تک‌تصویری، شبکهٔ سبک، پیمانه‌سازی کانولوشنی، توجه فضایی، بازآرایی زیرپیکسلی.

\chapter{مقدمه و صورت مسئله}

\section{تعریف مسئله}

در ابرتفکیک‌پذیری تک‌تصویری، ورودی یک تصویر کم‌وضوح
$x\in\mathbb{R}^{H\times W\times3}$ است و هدف یادگیری نگاشتی مانند
$f_{\theta}$ برای تولید تصویر
$\hat{y}=f_{\theta}(x)\in\mathbb{R}^{sH\times sW\times3}$ است؛ $s$ ضریب بزرگ‌نمایی و $\theta$ پارامترهای مدل است. چون در فرایند کاهش وضوح بخشی از بسامدهای بالا از بین می‌روند، پاسخ مسئله یکتا نیست. شبکه باید با استفاده از آمار دادهٔ آموزشی، لبه، بافت و ساختار محتمل را بازسازی کند.

روش‌های مبتنی بر شبکهٔ کانولوشنی در این مسئله موفق بوده‌اند، ولی وابستگی‌های دوربرد را به‌صورت غیرمستقیم و با افزایش عمق یا میدان دید می‌آموزند. ترنسفورمرها این محدودیت را با خودتوجهی کاهش می‌دهند، اما ماتریس توجه برای $N=HW$ نشانه، هزینه‌ای تقریباً متناسب با $N^2$ دارد. این هزینه روی تصویرهای بزرگ، مخصوصاً در مرحلهٔ استنتاج، مشکل اصلی است. مقالهٔ مبنا برای حفظ سبک ترنسفورمر و حذف ماتریس توجه، از پیمانه‌سازی کانولوشنی استفاده می‌کند \cite{Gendy2023,Hou2022}.

\section{هدف و دامنهٔ پروژه}

هدف ما فقط نوشتن لایهٔ مدل نبود. کل مسیر آزمایش را در نظر گرفتیم: ساخت دادهٔ کم‌وضوح با کاهش دو‌مکعبی، برش هم‌تراز وصله‌ها، افزایش داده، آموزش، ارزیابی، ذخیرهٔ مدل، تولید شکل و اجرای نسخه‌های حذف مؤلفه. سه ضریب مقیاس ۲، ۳ و ۴ در مدل پشتیبانی می‌شوند. پیکربندی پیش‌فرض مطابق آزمایش مقیاس دو مقاله است و حالت \lr{--quick} همان مدل کامل را روی داده‌ای کوچک اجرا می‌کند تا دستیار آموزشی بدون تغییر کد، در زمان کوتاه همهٔ مسیر را بررسی کند.

\section{جمع‌بندی مقالهٔ مبنا}

مقاله سه ادعای اصلی دارد. نخست، بلوک پیشنهادی به‌جای محاسبهٔ شباهت همهٔ مکان‌های تصویر، با کانولوشن و ضرب هادامارد کار می‌کند و نسبت به تعداد پیکسل‌ها رشد خطی دارد. دوم، قرار دادن این عملگر در ساختار پیش‌نرمال‌سازی‌شده و مانده، مزیت سازمان‌دهی بلوک ترنسفورمر را نگه می‌دارد. سوم، ترکیب چهار گروه از این بلوک‌ها با توجه فضایی تقویت‌شده و بازسازی زیرپیکسلی، نسبت مناسبی میان کیفیت، تعداد پارامتر و زمان اجرا ایجاد می‌کند. نویسندگان گزارش کرده‌اند که مدل آن‌ها نسبت به \lr{LWSwinIR} حدود $5.6$ برابر سریع‌تر است، در حالی که تعداد پارامتر و کیفیت دو مدل نزدیک‌اند \cite{Gendy2023}.

\chapter{روش پیشنهادی مقاله}

\section{پیمانه‌سازی کانولوشنی}

ورودی نرمال‌شده را با $F\in\mathbb{R}^{B\times C\times H\times W}$ نشان می‌دهیم. عملگر پیمانه‌سازی دو شاخه دارد. شاخهٔ اول زمینهٔ مکانی را با یک کانولوشن نقطه‌ای و سپس کانولوشن عمقی بزرگ استخراج می‌کند:

\begin{equation}
F_{d}=\operatorname{DWConv}_{k\times k}\!\left(\operatorname{Conv}_{1\times1}(F)\right).
\end{equation}

شاخهٔ دوم مقدار ویژگی را می‌سازد:

\begin{equation}
F_{v}=\operatorname{Conv}_{1\times1}(F).
\end{equation}

دو خروجی عضو‌به‌عضو ضرب و سپس فرافکنی می‌شوند:

\begin{equation}
\operatorname{CM}(F)=\operatorname{Conv}_{3\times3}(F_{d}\odot F_{v}).
\label{eq:modulation}
\end{equation}

تفاوت مهم نسخهٔ ابرتفکیک‌پذیری با عملگر اصلی \lr{Conv2Former} در کانولوشن آخر است. نسخهٔ اصلی از کانولوشن $1\times1$ استفاده می‌کرد، ولی مقاله آن را با $3\times3$ جایگزین کرده تا پیش از بازسازی، ویژگی‌های محلی بهتر حفظ شوند. نتیجهٔ حذف مؤلفه نیز این انتخاب را پشتیبانی می‌کند: در \lr{Urban100}، PSNR مدل کامل $32.58$ و نسخهٔ فرافکنی $1\times1$ برابر $32.48$ دسی‌بل گزارش شده است.

اگر هزینهٔ کانولوشن‌های نقطه‌ای را هم لحاظ کنیم، مرتبهٔ محاسباتی رابطهٔ \ref{eq:modulation} تقریباً برابر است با

\begin{equation}
\mathcal{O}\!\left(HWC^2 + HWk^2C\right),
\end{equation}

در حالی که خودتوجهی سراسری جمله‌ای از مرتبهٔ
$\mathcal{O}((HW)^2C)$ دارد. بنابراین با ثابت ماندن تعداد کانال و اندازهٔ کرنل، هزینهٔ پیمانه‌سازی نسبت به تعداد پیکسل‌ها خطی است.

\section{بلوک پیمانه‌سازی ترنسفورمرگونه}

هر بلوک از دو زیرلایه تشکیل می‌شود. ابتدا نرمال‌سازی لایه‌ای روی کانال‌ها انجام می‌شود، سپس پیمانه‌سازی کانولوشنی و اتصال مانده اعمال می‌شوند:

\begin{align}
Z &= F_{i-1} + \gamma_1\odot
\operatorname{CM}\!\left(\operatorname{LN}(F_{i-1})\right),\\
F_i &= Z + \gamma_2\odot
\operatorname{MLP}\!\left(\operatorname{LN}(Z)\right).
\label{eq:block}
\end{align}

$\gamma_1$ و $\gamma_2$ بردارهای مقیاس یادگرفتنی هستند. پرسپترون چندلایه ابتدا کانال‌ها را دو برابر می‌کند، یک کانولوشن عمقی $3\times3$ برای اختلاط مکانی دارد و سپس تعداد کانال‌ها را به مقدار اول برمی‌گرداند. تابع فعال‌سازی در این بخش \lr{GELU} است. وجود مسیرهای مانده باعث می‌شود گرادیان در شبکهٔ عمیق‌تر پایدارتر منتقل شود.

\section{گروه پیمانه‌سازی و توجه فضایی}

هر گروه از چهار بلوک رابطهٔ \ref{eq:block}، یک کانولوشن $3\times3$ و یک واحد توجه فضایی تقویت‌شده تشکیل شده است:

\begin{align}
F_i &= \operatorname{C2FB}_{i}(F_{i-1}),\qquad i=1,\ldots,n,\\
F_c &= \operatorname{Conv}_{3\times3}(F_n),\\
F_{\mathrm{group}} &= \operatorname{ESA}(F_c).
\end{align}

در واحد توجه فضایی، تعداد کانال‌ها از ۵۰ به ۱۲ کاهش می‌یابد. یک مسیر کم‌وضوح با کانولوشن گام‌دار و بیشینه‌گیری، میدان دید را بزرگ می‌کند و یک مسیر کوتاه اطلاعات محلی را نگه می‌دارد. پس از درون‌یابی مسیر کم‌وضوح و جمع دو مسیر، یک نگاشت سیگموئیدی به‌عنوان ماسک توجه روی ورودی ضرب می‌شود. این واحد از ایدهٔ توجه فضایی تقویت‌شده در شبکهٔ تجمیع ویژگی مانده استفاده می‌کند \cite{Liu2020}.

\section{چارچوب کامل STSN}

مدل کامل سه بخش دارد. در بخش اول، یک کانولوشن $3\times3$ ویژگی کم‌عمق را استخراج می‌کند:

\begin{equation}
F_0=\operatorname{Conv}_{3\times3}(x).
\end{equation}

در بخش دوم، چهار گروه به‌صورت زنجیره‌ای اجرا می‌شوند:

\begin{equation}
F_g=\operatorname{Group}_g(F_{g-1}),\qquad g=1,\ldots,4.
\end{equation}

خروجی کم‌عمق و خروجی همهٔ گروه‌ها روی بعد کانال به هم چسبانده می‌شوند. یک کانولوشن $1\times1$ تعداد کانال‌ها را کاهش می‌دهد و کانولوشن $3\times3$ بعدی ویژگی‌ها را هموار و مخلوط می‌کند:

\begin{equation}
F_{\mathrm{comb}}=\operatorname{Conv}_{3\times3}\!\left(
\operatorname{Conv}_{1\times1}\!\left([F_0,F_1,F_2,F_3,F_4]\right)\right).
\end{equation}

در پایان، $F_{\mathrm{comb}}$ با $F_0$ جمع می‌شود. کانولوشن بازسازی $3s^2$ کانال تولید می‌کند و بازآرایی زیرپیکسلی آن‌ها را به تصویر RGB با ابعاد $sH\times sW$ تبدیل می‌کند \cite{Shi2016}. شکل \ref{fig:architecture} مسیر دادهٔ پیاده‌سازی ما را نشان می‌دهد. این شکل بازترسیم مستقل معماری است و تصویر صفحه یا کپی شکل مقاله نیست.

\begin{figure}[H]
    \centering
    \includegraphics[width=\textwidth]{images/stsn_architecture.png}
    \caption{مسیر داده در پیاده‌سازی STSN؛ چهار گروه میانی هرکدام شامل چهار بلوک پیمانه‌سازی هستند. بازترسیم توسط ما بر اساس روابط مقاله.}
    \label{fig:architecture}
\end{figure}

\chapter{جزئیات پیاده‌سازی}

\section{ساختار منبع کد}

کد در پوشهٔ \lr{homework\_code} قرار دارد و ساختار درختی خواسته‌شده رعایت شده است. فایل‌های اضافی \lr{ablation.py}، \lr{download\_set5.py} و پوشهٔ آزمون‌ها برای بازتولیدپذیری افزوده شده‌اند.

\begin{latin}
\begin{verbatim}
homework_code/
|-- data/
|   |-- data_loader.py
|   |-- download_set5.py
|   `-- dataset/
|-- models/
|   |-- model.py
|   `-- saved_models/
|-- scripts/
|   |-- main.py
|   |-- train.py
|   |-- evaluate.py
|   `-- ablation.py
|-- config/
|   |-- config.yaml
|   `-- logging.yaml
|-- utils/
|   |-- config.py
|   |-- visualization.py
|   `-- metrics.py
|-- tests/test_pipeline.py
|-- README.md
`-- requirements.txt
\end{verbatim}
\end{latin}

همهٔ مسیرهای موجود در فایل پیکربندی نسبت به ریشهٔ پروژه تفسیر می‌شوند. تنها مسیر مطلقی که در زمان اجرا دیده می‌شود با \lr{Path(\_\_file\_\_).resolve()} از محل خود فایل ساخته می‌شود و هیچ نام کاربری، حرف درایو یا مسیر وابسته به رایانه در کد نوشته نشده است.

\section{استنتاج جزئیات نامشخص معماری}

متن مقاله نسبت گسترش پرسپترون، اندازهٔ کرنل عمقی و تعداد کانال میانی توجه فضایی را صریح ذکر نکرده است. به‌جای انتخاب دلخواه، از جدول‌های حذف مؤلفه به‌عنوان قید استفاده کردیم. منطق این بازسازی به‌صورت زیر بود:

\begin{itemize}
    \item تفاوت مدل کامل و مدل بدون پرسپترون $180{,}800$ پارامتر است. با ۱۶ بلوک، سهم هر پرسپترون همراه با نرمال‌سازی و مقیاس لایه $11{,}300$ می‌شود؛ این عدد دقیقاً با نسبت گسترش ۲ و کانولوشن عمقی $3\times3$ سازگار است.
    \item تفاوت مدل کامل و مدل بدون پیمانه‌سازی $542{,}400$ پارامتر است. سهم $33{,}900$ پارامتری هر بلوک، با کرنل عمقی $11\times11$، دو کانولوشن نقطه‌ای، فرافکنی $3\times3$، نرمال‌سازی و مقیاس لایه دقیقاً به‌دست می‌آید.
    \item حذف توجه فضایی $26{,}600$ پارامتر کم می‌کند. تقسیم این مقدار میان چهار گروه، $6{,}650$ پارامتر برای هر واحد می‌دهد که با ۱۲ کانال میانی، یعنی $\lfloor50/4\rfloor$، برابر است.
\end{itemize}

این بررسی یک آزمون خودکار شده است؛ بنابراین تغییر ناخواستهٔ معماری باعث شکست آزمون می‌شود. جدول \ref{tab:param-ablation} نتیجه را نشان می‌دهد.

\begin{table}[H]
\centering
\caption{تطبیق شمار پارامترهای پیاده‌سازی با مقادیر جدول‌های حذف مؤلفهٔ مقاله در مقیاس دو}
\label{tab:param-ablation}
\begin{tabular}{lrrr}
\toprule
گونهٔ مدل & مقاله (هزار) & پیاده‌سازی (دقیق) & اختلاف \\
\midrule
مدل کامل & ۸۸۱٫۹ & ۸۸۱٬۹۱۲ & ۰ \\
فرافکنی $1\times1$ & ۵۶۱٫۹ & ۵۶۱٬۹۱۲ & ۰ \\
بدون دروازهٔ توجه & ۷۰۲٫۷ & ۷۰۲٬۷۱۲ & ۰ \\
بدون پیمانه‌سازی & ۳۳۹٫۵ & ۳۳۹٬۵۱۲ & ۰ \\
بدون پرسپترون & ۷۰۱٫۱ & ۷۰۱٬۱۱۲ & ۰ \\
بدون توجه فضایی & ۸۵۵٫۳ & ۸۵۵٬۳۱۲ & ۰ \\
بدون کانولوشن گروه & ۷۹۱٫۷ & ۷۹۱٬۷۱۲ & ۰ \\
\bottomrule
\end{tabular}
\end{table}

منظور از اختلاف صفر در جدول، برابری با مقدار دقیق استنتاج‌شده پیش از گرد کردن به هزار پارامتر است؛ خود مقاله اعداد ستون دوم را با یک رقم اعشار گزارش کرده است.

\section{آماده‌سازی داده}

کلاس مجموعه‌داده دو حالت دارد. اگر جفت LR موجود باشد، نام فایل‌های دو پوشه تطبیق داده و ابعاد آن‌ها هم‌تراز می‌شود. در غیر این صورت، تصویر HR ابتدا طوری برش می‌خورد که ابعادش مضرب $s$ باشد و سپس کاهش دو‌مکعبی روی آن اجرا می‌شود. در حالت آموزش یک وصلهٔ تصادفی HR انتخاب و مختصات LR متناظر از مضرب دقیق آن برداشته می‌شود؛ در نتیجه خطای نیم‌پیکسل یا جابه‌جایی میان هدف و ورودی ایجاد نمی‌شود.

افزایش داده شامل بازتاب افقی تصادفی و چرخش تصادفی صفر، ۹۰، ۱۸۰ یا ۲۷۰ درجه است. این تنظیم دقیقاً با توضیح مقاله سازگار است. داده‌ها به تنسور RGB در بازهٔ صفر تا یک تبدیل می‌شوند. برای اجرای آفلاین، تابعی نوشته‌ایم که تصویرهای ساختگی دارای گرادیان، بافت تناوبی، لبه، دایره و مستطیل تولید می‌کند. این تصاویر از منبع خارجی گرفته نشده‌اند.

در اجرای نهایی، هر ۸۰۰ تصویر آموزشی \lr{DIV2K} و جفت‌های کم‌وضوح رسمی دو‌مکعبی برای هر سه مقیاس استفاده شدند \cite{Agustsson2017}. ارزیابی نیز روی تمام تصاویر پنج مجموعهٔ \lr{Set5}، \lr{Set14}، \lr{B100}، \lr{Urban100} و \lr{Manga109} انجام شد. برای جلوگیری از اختلاف میان پیاده‌سازی‌های بازنمونه‌برداری، جفت‌های آمادهٔ \lr{Set5} و \lr{Set14} مستقیماً به‌کار رفتند و برای سه مجموعهٔ دیگر، تصاویر کم‌وضوح با تابع \lr{imresize} متلب و گزینهٔ دو‌مکعبیِ پادهمپوشانی ساخته شدند. تعداد تصاویر پیش از شروع آموزش و ارزیابی به‌صورت برنامه‌ای کنترل شد.

\section{آموزش و ذخیرهٔ مدل}

تابع هدف میانگین قدرمطلق خطای پیکسل‌ها است:

\begin{equation}
\mathcal{L}_1(\theta)=\frac{1}{3sHsW}\sum_{c=1}^{3}\sum_{u=1}^{sH}\sum_{v=1}^{sW}
\left|f_{\theta}(x)_{cuv}-y_{cuv}\right|.
\end{equation}

بهینه‌ساز آدم با $\beta_1=0.9$، $\beta_2=0.99$ و $\epsilon=10^{-8}$ استفاده شده است \cite{Kingma2014}. نرخ یادگیری اولیه $5\times10^{-4}$ است و هر ۲۰۰ دوره نصف می‌شود. مقاله اندازهٔ دستهٔ ۳۲، آموزش ۱۰۰۰‌دوره‌ای و سپس بازآموزی با راهبرد شروع گرم را گزارش کرده است. مطابق تصمیم پروژه، ما همان اندازهٔ دسته و همهٔ ۸۰۰ تصویر را حفظ کردیم، اما بودجهٔ هر مرحله را به ۱۰۰ دوره کاهش دادیم: مرحلهٔ نخست از وزن تصادفی و مرحلهٔ دوم از وزن پایان مرحلهٔ نخست آغاز می‌شود. در هر مرحله بهینه‌ساز و زمان‌بند از نو ساخته می‌شوند؛ ازاین‌رو نرخ یادگیری در بازهٔ ۱۰۰‌دوره‌ای هنوز به نخستین پلهٔ کاهشِ دورهٔ ۲۰۰ نرسیده و ثابت می‌ماند. نقاط وارسی آخرین دوره و بهترین PSNR پایشی ذخیره می‌شوند، ولی اعداد اصلی از نقطهٔ وارسی پایان دورهٔ ۱۰۰ مرحلهٔ شروع گرم و اعداد «پیش از شروع گرم» از پایان مرحلهٔ نخست گرفته می‌شوند؛ بنابراین Set5 در انتخاب وزن گزارش‌شده دخالت ندارد.

\section{سنجه‌های ارزیابی}

مطابق مقاله، ارزیابی روی کانال روشنایی انجام می‌شود. تبدیل RGB به Y با قرارداد رایج MATLAB به‌شکل زیر است:

\begin{equation}
Y=\frac{65.481R+128.553G+24.966B+16}{255},
\end{equation}

که در آن $R$، $G$ و $B$ در بازهٔ صفر تا یک‌اند. سپس حاشیه‌ای به عرض $s$ پیکسل حذف می‌شود. نسبت بیشینهٔ سیگنال به نویز از خطای میانگین مربعات محاسبه می‌شود:

\begin{equation}
\operatorname{PSNR}=10\log_{10}\left(\frac{1}{\operatorname{MSE}}\right).
\end{equation}

شاخص شباهت ساختاری با پنجرهٔ گاوسی ۱۱ پیکسلی، انحراف معیار $1.5$ و بازهٔ دادهٔ یک محاسبه می‌شود \cite{Wang2004}. برای مقایسهٔ منصفانه، همین تبدیل رنگ و حذف حاشیه هم برای STSN و هم برای خط پایهٔ دو‌مکعبی استفاده شده است.

\chapter{آزمایش‌ها و نتایج}

\section{تنظیمات گزارش‌شده در مقاله}

نویسندگان برای ضرایب ۲، ۳ و ۴ به‌ترتیب اندازهٔ وصلهٔ ۹۶، ۱۴۴ و ۱۹۲ را به‌کار برده‌اند. چهار گروه، چهار بلوک در هر گروه و ۵۰ کانال ویژگی ثابت مانده‌اند. آموزش روی \lr{DIV2K} و ارزیابی روی \lr{Set5}، \lr{Set14}، \lr{B100}، \lr{Urban100} و \lr{Manga109} انجام شده است. جدول \ref{tab:paper-main-results} فقط ردیف STSN جدول اصلی مقاله را بازنویسی می‌کند تا مرجع عددی روش روشن باشد؛ این اعداد خروجی اجرای ما نیستند.

\begin{table}[H]
\centering
\caption{نتایج STSN گزارش‌شده در مقاله؛ هر خانه به‌ترتیب PSNR برحسب دسی‌بل و SSIM را نشان می‌دهد}
\label{tab:paper-main-results}
\resizebox{\textwidth}{!}{%
\begin{tabular}{c c c c c c c}
\toprule
مقیاس & پارامتر (هزار) & Set5 & Set14 & B100 & Urban100 & Manga109 \\
\midrule
$\times2$ & ۸۸۱٫۹ & ۳۸٫۱۹ / ۰٫۹۶۱۱ & ۳۳٫۷۸ / ۰٫۹۱۹۹ & ۳۲٫۳۰ / ۰٫۹۰۱۳ & ۳۲٫۶۸ / ۰٫۹۳۳۶ & ۳۹٫۱۳ / ۰٫۹۷۷۸ \\
$\times3$ & ۸۸۸٫۷ & ۳۴٫۶۲ / ۰٫۹۲۹۲ & ۳۰٫۵۴ / ۰٫۸۴۶۶ & ۲۹٫۲۲ / ۰٫۸۰۹۰ & ۲۸٫۵۹ / ۰٫۸۶۲۱ & ۳۴٫۱۱ / ۰٫۹۴۸۰ \\
$\times4$ & ۸۹۸٫۲ & ۳۲٫۴۶ / ۰٫۸۹۸۲ & ۲۸٫۷۶ / ۰٫۷۸۶۰ & ۲۷٫۶۸ / ۰٫۷۴۰۵ & ۲۶٫۳۹ / ۰٫۷۹۷۱ & ۳۰٫۹۳ / ۰٫۹۱۴۲ \\
\bottomrule
\end{tabular}}
\end{table}

زمان گزارش‌شده برای مقیاس‌های ۲، ۳ و ۴ به‌ترتیب ۶۴۰، ۲۹۸ و ۱۶۸ میلی‌ثانیه و تعداد ضرب‌ـ‌جمع‌ها به‌ترتیب ۱۹۷٫۷، ۹۹٫۹ و ۵۰٫۳ میلیارد است. کاهش هزینه با افزایش ضریب مقیاس در نگاه اول عجیب به نظر می‌رسد، ولی ابعاد ورودی LR برای خروجی ثابت کوچک‌تر می‌شود. در مقیاس دو، زمان ۶۴۰ میلی‌ثانیه در برابر ۳۵۹۰ میلی‌ثانیهٔ \lr{LWSwinIR} همان نسبت سرعت حدود $5.6$ برابر را می‌سازد.

\section{آزمایش حذف مؤلفهٔ مقاله}

مهم‌ترین نتیجهٔ حذف مؤلفه روی مجموعه‌های دارای ساختار زیاد دیده می‌شود. جدول \ref{tab:paper-ablation-quality} تغییر PSNR دو مجموعه را نشان می‌دهد. حذف کامل پیمانه‌سازی در \lr{Urban100} حدود $0.52$ دسی‌بل و در \lr{Manga109} حدود $0.36$ دسی‌بل افت ایجاد کرده است. حذف پرسپترون نیز افت قابل‌توجه دارد. توجه فضایی و کانولوشن محلی گروه پارامتر بسیار کمتری نسبت به بدنه دارند، ولی هر دو روی تصاویر شهری اثر محسوسی گذاشته‌اند.

\begin{table}[H]
\centering
\caption{بخشی از نتایج کیفی حذف مؤلفهٔ مقاله پیش از شروع گرم در مقیاس دو}
\label{tab:paper-ablation-quality}
\begin{tabular}{lcc}
\toprule
گونهٔ مدل & PSNR در Urban100 & PSNR در Manga109 \\
\midrule
مدل کامل & ۳۲٫۵۸ & ۳۹٫۰۶ \\
فرافکنی $1\times1$ & ۳۲٫۴۸ & ۳۹٫۰۱ \\
بدون دروازهٔ توجه & ۳۲٫۲۰ & ۳۸٫۸۵ \\
بدون پیمانه‌سازی & ۳۲٫۰۶ & ۳۸٫۷۰ \\
بدون پرسپترون & ۳۲٫۳۰ & ۳۸٫۸۶ \\
بدون توجه فضایی & ۳۲٫۴۳ & ۳۹٫۰۰ \\
بدون کانولوشن گروه & ۳۲٫۴۴ & ۳۹٫۰۰ \\
\bottomrule
\end{tabular}
\end{table}

شروع گرم بدون افزودن پارامتر، PSNR را روی \lr{Urban100} از ۳۲٫۵۸ به ۳۲٫۶۸ و روی \lr{Manga109} از ۳۹٫۰۶ به ۳۹٫۱۳ رسانده است. به همین دلیل کد ما مسیر شروع گرم را جدا از ادامهٔ عادی آموزش پیاده می‌کند.

\section{تنظیمات اجرای نهایی ما}

اجرای نهایی روی کارت گرافیکی \lr{NVIDIA Tesla T4} در محیط کگل، پایتون ۳٫۱۲٫۱۳، \lr{PyTorch 2.10.0+cu128} و کودا ۱۲٫۸ انجام شد. برای هر سه مقیاس، همهٔ ۸۰۰ تصویر آموزشی \lr{DIV2K} و جفت‌های رسمی کم‌وضوح دو‌مکعبی به‌کار رفتند. اندازهٔ وصلهٔ پروضوح برای مقیاس‌های ۲، ۳ و ۴ به‌ترتیب ۹۶، ۱۴۴ و ۱۹۲ و اندازهٔ وصلهٔ کم‌وضوح در هر سه حالت ۴۸ بود. اندازهٔ دسته ۳۲، تعداد کارگر بارگذاری داده ۲، بذر تصادفی ۸۱۰۱۰۴۰۳۹ و زیان همان $L_1$ مقاله بود.

برای هر مقیاس دو مرحله اجرا کردیم: مرحلهٔ نخست ۱۰۰ دوره از وزن تصادفی و مرحلهٔ دوم ۱۰۰ دوره با شروع از وزن پایان مرحلهٔ نخست. هر مرحله ۲۵۰۰ به‌روزرسانی و کل آموزش هر مقیاس ۵۰۰۰ به‌روزرسانی داشت. مقاله برای آموزش نخست ۱۰۰۰ دوره را مشخص کرده، اما تعداد دورهٔ بازآموزی مدل اصلی را صریح ننوشته است؛ بنابراین فرض مستند ما این بود که بودجهٔ ۱۰۰‌دوره‌ای انتخاب‌شده برای پروژه، به هر یک از دو مرحله اختصاص یابد. برای جلوگیری از انتخاب نتیجه بر پایهٔ مجموعهٔ آزمون، تمام اعداد اصلی از نقطهٔ وارسی آخرین دوره گرفته شدند، نه بهترین نقطهٔ وارسی بر اساس \lr{Set5}.

جدول \ref{tab:our-training-time} زمان واقعی آموزش را نشان می‌دهد. مجموع زمان پردازش کارت گرافیکی شش اجرا حدود ۹٫۴۱ ساعت بود. اجراهای سه مقیاس مستقل از یکدیگرند و به همین علت امکان اجرای هم‌زمان آن‌ها وجود داشت.

\begin{table}[H]
\centering
\caption{زمان و شمار پارامتر اجرای نهایی؛ زمان‌ها از تاریخچهٔ ذخیره‌شدهٔ خود برنامه استخراج شده‌اند}
\label{tab:our-training-time}
\begin{tabular}{rrrrr}
\toprule
مقیاس & شمار پارامتر & مرحلهٔ نخست (دقیقه) & شروع گرم (دقیقه) & مجموع (دقیقه) \\
\midrule
۲ & ۸۸۱٬۹۱۲ & ۹۵٫۱ & ۹۴٫۴ & ۱۸۹٫۵ \\
۳ & ۸۸۸٬۶۷۷ & ۸۷٫۸ & ۹۶٫۸ & ۱۸۴٫۶ \\
۴ & ۸۹۸٬۱۴۸ & ۹۰٫۷ & ۹۹٫۶ & ۱۹۰٫۳ \\
\bottomrule
\end{tabular}
\end{table}

\section{روند همگرایی}

جدول \ref{tab:our-convergence} زیان نخستین و آخرین دورهٔ هر مرحله را مقایسه می‌کند. در هر سه مقیاس، زیان پایان مرحلهٔ شروع گرم از پایان آموزش نخست کمتر شد. بزرگ‌تر بودن زیان دورهٔ نخست شروع گرم نسبت به آخرین دورهٔ مرحلهٔ نخست تناقض نیست، زیرا ترتیب وصله‌های تصادفی، بهینه‌ساز و حالت‌های مومنتوم در ابتدای مرحلهٔ دوم از نو ساخته شدند.

\begin{table}[H]
\centering
\caption{تغییر میانگین زیان $L_1$ روی دادهٔ کامل \lr{DIV2K}}
\label{tab:our-convergence}
\begin{tabular}{rrrrr}
\toprule
مقیاس & نخستین دورهٔ مرحلهٔ اول & پایان مرحلهٔ اول & نخستین دورهٔ شروع گرم & پایان شروع گرم \\
\midrule
۲ & ۰٫۱۵۷۳۳ & ۰٫۰۱۵۴۲ & ۰٫۰۲۰۹۷ & ۰٫۰۱۴۵۳ \\
۳ & ۰٫۱۹۹۴۳ & ۰٫۰۲۳۳۰ & ۰٫۰۲۹۶۰ & ۰٫۰۲۲۶۵ \\
۴ & ۰٫۱۶۳۵۸ & ۰٫۰۲۸۶۲ & ۰٫۰۳۵۳۶ & ۰٫۰۲۷۰۰ \\
\bottomrule
\end{tabular}
\end{table}

شکل‌های \ref{fig:our-x2-curves} تا \ref{fig:our-x4-curves} هر شش نمودار آموزش را نمایش می‌دهند؛ برای هر مقیاس، نمودار مرحلهٔ نخست و شروع گرم جداگانه آمده است. هر نقطهٔ آبی میانگین زیان تمام ۲۵ دستهٔ آن دوره و نقاط قرمز PSNR پایشی روی \lr{Set5} هستند. افت ناگهانی یا پلهٔ نرخ یادگیری دیده نمی‌شود، زیرا نخستین نصف‌شدن نرخ یادگیری در دورهٔ ۲۰۰ تعریف شده و هر مرحلهٔ ما در دورهٔ ۱۰۰ پایان می‌یابد.

\begin{figure}[H]
    \centering
    \includegraphics[width=0.485\textwidth]{images/training_curve_x2_stage1.png}
    \hfill
    \includegraphics[width=0.485\textwidth]{images/training_curve_x2_warm.png}
    \caption{روند آموزش مدل مقیاس دو روی همهٔ ۸۰۰ تصویر \lr{DIV2K}؛ سمت راست مرحلهٔ ۱۰۰‌دوره‌ای از صفر و سمت چپ مرحلهٔ ۱۰۰‌دوره‌ای شروع گرم است. محور و مقادیر مستقیماً به‌وسیلهٔ برنامه تولید شده‌اند.}
    \label{fig:our-x2-curves}
\end{figure}

\begin{figure}[H]
    \centering
    \includegraphics[width=0.485\textwidth]{images/training_curve_x3_stage1.png}
    \hfill
    \includegraphics[width=0.485\textwidth]{images/training_curve_x3_warm.png}
    \caption{روند کامل آموزش مقیاس سه؛ سمت راست آموزش ۱۰۰‌دوره‌ای از صفر و سمت چپ آموزش ۱۰۰‌دوره‌ای شروع گرم است. هر دو نمودار از تاریخچهٔ ذخیره‌شدهٔ اجرای کگل رسم شده‌اند.}
    \label{fig:our-x3-curves}
\end{figure}

\begin{figure}[H]
    \centering
    \includegraphics[width=0.485\textwidth]{images/training_curve_x4_stage1.png}
    \hfill
    \includegraphics[width=0.485\textwidth]{images/training_curve_x4_warm.png}
    \caption{روند کامل آموزش مقیاس چهار؛ سمت راست آموزش ۱۰۰‌دوره‌ای از صفر و سمت چپ آموزش ۱۰۰‌دوره‌ای شروع گرم است. نوسان PSNR پایشی در این مقیاس علت گزارش‌کردن وزن آخرین دوره به‌جای انتخاب بهترین وزن را روشن می‌کند.}
    \label{fig:our-x4-curves}
\end{figure}

در مقیاس دو، زیان شروع گرم در پایان به ۰٫۰۱۴۵۳ رسید و بیشترین PSNR پایشی نیز در همان انتهای مرحله به‌دست آمد. مقیاس سه روندی مشابه داشت و زیان نهایی آن ۰٫۰۲۲۶۵ شد. در مقیاس چهار، زیان تا ۰٫۰۲۷۰۰ کاهش یافت، اما PSNR پایشی به‌علت حساسیت بیشتر بازسازی چهاربرابری نوسان داشت؛ بهترین مقدار پایشی ۳۰٫۹۶۷ و مقدار آخرین دوره ۳۰٫۷۷۹ دسی‌بل بود. ما مقدار دوم را در تمام جدول‌های اصلی به‌کار بردیم تا قاعدهٔ انتخاب وزن برای هر سه مقیاس یکسان بماند.

شکل \ref{fig:combined-training-loss} همین روند را به‌صورت فشرده نمایش می‌دهد. در هر سه مقیاس، منحنی مرحلهٔ شروع گرم از مقدار زیان پایین‌تری آغاز می‌شود و در پایان نیز از مرحلهٔ نخست کمتر است.

\begin{figure}[H]
    \centering
    \includegraphics[width=\textwidth]{images/training_loss_curves.png}
    \caption{مقایسهٔ زیان آموزش در دو مرحلهٔ ۱۰۰‌دوره‌ای برای هر سه مقیاس. منحنی‌ها مستقیماً از تاریخچه‌های ذخیره‌شدهٔ اجرای نهایی رسم شده‌اند.}
    \label{fig:combined-training-loss}
\end{figure}

\section{نتایج نهایی روی پنج مجموعهٔ معیار}

برای اینکه اثر هر مرحله مستقل دیده شود، جدول \ref{tab:our-stage-one-results} همهٔ نتایج پایان آموزش نخست را گزارش می‌کند. سپس جدول \ref{tab:our-full-results} نتیجهٔ پایان دورهٔ ۱۰۰ شروع گرم را برای همان ۱۵ ترکیب مقیاس و مجموعه می‌آورد. سنجه‌ها روی کانال روشنایی و پس از حذف حاشیه‌ای به عرض ضریب مقیاس محاسبه شده‌اند. هر مقدار میانگین همهٔ تصاویر آن مجموعه است، نه نتیجهٔ یک نمونهٔ منتخب: ۵ تصویر \lr{Set5}، ۱۴ تصویر \lr{Set14}، ۱۰۰ تصویر \lr{B100}، ۱۰۰ تصویر \lr{Urban100} و ۱۰۹ تصویر \lr{Manga109}؛ در مجموع ۳۲۸ تصویر در هر اجرا. نتیجهٔ تصویر‌به‌تصویر شش اجرا نیز در ۳۰ فایل CSV داخل پوشه‌های \lr{outputs} تحویل داده شده است. این ۱۹۶۸ ردیف خام مبنای میانگین‌های گزارش‌اند، اما برای خوانایی به‌صورت کامل داخل متن چاپ نشده‌اند.

\begin{table}[H]
\centering
\caption{همهٔ نتایج پایان مرحلهٔ نخست ۱۰۰‌دوره‌ای پیش از شروع گرم؛ هر خانه به‌ترتیب PSNR برحسب دسی‌بل و SSIM را نشان می‌دهد}
\label{tab:our-stage-one-results}
\resizebox{\textwidth}{!}{%
\begin{tabular}{rrrrrr}
\toprule
مقیاس & \lr{Set5} & \lr{Set14} & \lr{B100} & \lr{Urban100} & \lr{Manga109} \\
\midrule
۲ & ۳۶٫۷۵۹ / ۰٫۹۵۴۳ & ۳۲٫۴۸۲ / ۰٫۹۰۶۶ & ۳۱٫۳۸۷ / ۰٫۸۸۸۰ & ۲۹٫۶۶۸ / ۰٫۸۹۷۳ & ۳۶٫۰۷۵ / ۰٫۹۶۵۳ \\
۳ & ۳۲٫۴۶۷ / ۰٫۹۰۲۳ & ۲۹٫۱۱۴ / ۰٫۸۱۶۵ & ۲۸٫۲۶۰ / ۰٫۷۸۲۹ & ۲۵٫۹۳۹ / ۰٫۷۹۱۷ & ۲۹٫۸۶۱ / ۰٫۸۹۹۰ \\
۴ & ۳۰٫۲۵۴ / ۰٫۸۵۷۳ & ۲۷٫۳۵۰ / ۰٫۷۵۰۸ & ۲۶٫۷۹۳ / ۰٫۷۱۱۰ & ۲۴٫۳۲۳ / ۰٫۷۱۸۴ & ۲۷٫۰۷۶ / ۰٫۸۴۳۵ \\
\bottomrule
\end{tabular}}
\end{table}

مرحلهٔ نخست در همهٔ مجموعه‌ها از دو‌مکعبی بهتر شد، ولی فاصلهٔ آن با خروجی شروع گرم روی مجموعه‌های بافت‌دار بیشتر است. برای نمونه، شروع گرم PSNR مقیاس سه را روی \lr{Manga109} از ۲۹٫۸۶۱ به ۳۱٫۲۶۹ رساند؛ یعنی افزایش ۱٫۴۰۸ دسی‌بل بدون تغییر معماری یا شمار پارامتر.

\begin{table}[H]
\centering
\caption{همهٔ نتایج پایان مرحلهٔ شروع گرم در مقایسه با درون‌یابی دو‌مکعبی و STSN گزارش‌شده در مقاله؛ هر جفت ستون شامل PSNR برحسب دسی‌بل و SSIM است}
\label{tab:our-full-results}
\resizebox{\textwidth}{!}{%
\begin{tabular}{clrrrrrr}
\toprule
مقیاس & مجموعه & \multicolumn{2}{c}{مدل ما} & \multicolumn{2}{c}{دو‌مکعبی} & \multicolumn{2}{c}{مقاله} \\
 & & PSNR & SSIM & PSNR & SSIM & PSNR & SSIM \\
\midrule
۲ & \lr{Set5}     & ۳۷٫۱۶۳ & ۰٫۹۵۷۵ & ۳۳٫۹۶۰ & ۰٫۹۳۳۶ & ۳۸٫۱۹ & ۰٫۹۶۱۱ \\
۲ & \lr{Set14}    & ۳۲٫۸۲۶ & ۰٫۹۱۱۵ & ۳۰٫۵۴۴ & ۰٫۸۷۵۳ & ۳۳٫۷۸ & ۰٫۹۱۹۹ \\
۲ & \lr{B100}     & ۳۱٫۶۸۵ & ۰٫۸۹۳۶ & ۲۹٫۷۴۳ & ۰٫۸۴۹۷ & ۳۲٫۳۰ & ۰٫۹۰۱۳ \\
۲ & \lr{Urban100} & ۳۰٫۳۴۳ & ۰٫۹۰۹۳ & ۲۷٫۰۷۳ & ۰٫۸۴۵۸ & ۳۲٫۶۸ & ۰٫۹۳۳۶ \\
۲ & \lr{Manga109} & ۳۶٫۸۶۷ & ۰٫۹۷۲۰ & ۳۱٫۲۵۲ & ۰٫۹۳۸۶ & ۳۹٫۱۳ & ۰٫۹۷۷۸ \\
\midrule
۳ & \lr{Set5}     & ۳۳٫۱۷۹ & ۰٫۹۱۶۶ & ۳۰٫۶۱۳ & ۰٫۸۷۲۱ & ۳۴٫۶۲ & ۰٫۹۲۹۲ \\
۳ & \lr{Set14}    & ۲۹٫۴۶۶ & ۰٫۸۲۷۷ & ۲۷٫۷۸۷ & ۰٫۷۸۱۲ & ۳۰٫۵۴ & ۰٫۸۴۶۶ \\
۳ & \lr{B100}     & ۲۸٫۵۷۰ & ۰٫۷۹۳۲ & ۲۷٫۳۱۹ & ۰٫۷۴۴۳ & ۲۹٫۲۲ & ۰٫۸۰۹۰ \\
۳ & \lr{Urban100} & ۲۶٫۶۵۵ & ۰٫۸۱۶۴ & ۲۴٫۵۸۲ & ۰٫۷۳۹۸ & ۲۸٫۵۹ & ۰٫۸۶۲۱ \\
۳ & \lr{Manga109} & ۳۱٫۲۶۹ & ۰٫۹۲۳۶ & ۲۷٫۱۹۸ & ۰٫۸۶۰۳ & ۳۴٫۱۱ & ۰٫۹۴۸۰ \\
\midrule
۴ & \lr{Set5}     & ۳۰٫۷۷۹ & ۰٫۸۷۲۲ & ۲۸٫۶۰۱ & ۰٫۸۱۴۲ & ۳۲٫۴۶ & ۰٫۸۹۸۲ \\
۴ & \lr{Set14}    & ۲۷٫۷۲۳ & ۰٫۷۶۲۴ & ۲۶٫۲۱۰ & ۰٫۷۰۹۰ & ۲۸٫۷۶ & ۰٫۷۸۶۰ \\
۴ & \lr{B100}     & ۲۷٫۰۵۸ & ۰٫۷۱۹۷ & ۲۶٫۰۵۱ & ۰٫۶۷۲۲ & ۲۷٫۶۸ & ۰٫۷۴۰۵ \\
۴ & \lr{Urban100} & ۲۴٫۷۸۱ & ۰٫۷۳۹۳ & ۲۳٫۲۴۰ & ۰٫۶۶۱۶ & ۲۶٫۳۹ & ۰٫۷۹۷۱ \\
۴ & \lr{Manga109} & ۲۸٫۱۶۸ & ۰٫۸۶۹۷ & ۲۵٫۰۷۲ & ۰٫۷۹۰۶ & ۳۰٫۹۳ & ۰٫۹۱۴۲ \\
\bottomrule
\end{tabular}}
\end{table}

شکل \ref{fig:benchmark-psnr} مقایسهٔ جدول را به‌صورت بصری خلاصه می‌کند. در همهٔ مجموعه‌ها و هر سه مقیاس، خروجی نهایی STSN از درون‌یابی دو‌مکعبی PSNR بالاتری دارد.

\begin{figure}[H]
    \centering
    \includegraphics[width=\textwidth]{images/benchmark_psnr_vs_bicubic.png}
    \caption{مقایسهٔ PSNR خروجی نهایی STSN پس از شروع گرم با خط پایهٔ دو‌مکعبی روی پنج مجموعهٔ معیار. هر پنل یک ضریب بزرگ‌نمایی را نشان می‌دهد.}
    \label{fig:benchmark-psnr}
\end{figure}

در همهٔ ۱۵ حالت، مدل آموزش‌دیده از درون‌یابی دو‌مکعبی بهتر است. بیشترین بهبود در \lr{Manga109} رخ داد: ۵٫۶۱۵، ۴٫۰۷۱ و ۳٫۰۹۶ دسی‌بل برای مقیاس‌های ۲، ۳ و ۴. این رفتار با ماهیت این مجموعه سازگار است، زیرا خطوط و بافت‌های تکراری تصویرهای کمیک از یادگیری ساختاری بیش از هموارسازی ثابت سود می‌برند. کمترین بهبود مربوط به \lr{B100} بود، ولی حتی آن‌جا نیز برای هر سه مقیاس بیش از یک دسی‌بل افزایش به‌دست آمد.

\subsection{تحلیل مجموعه‌به‌مجموعه}

در \lr{Set5} که کوچک‌ترین و معمولاً ساده‌ترین معیار است، خروجی نهایی به‌ترتیب ۳۷٫۱۶۳، ۳۳٫۱۷۹ و ۳۰٫۷۷۹ دسی‌بل شد. SSIM نیز از ۰٫۹۵۷۵ در مقیاس دو به ۰٫۸۷۲۲ در مقیاس چهار کاهش یافت. این افت با افزایش ضریب بزرگ‌نمایی طبیعی است، زیرا در مقیاس چهار تعداد بیشتری از جزئیات باید از ورودی کوچک‌تر تخمین زده شوند.

در \lr{Set14} تنوع محتوایی بیشتر است و نتیجهٔ سه مقیاس به ۳۲٫۸۲۶، ۲۹٫۴۶۶ و ۲۷٫۷۲۳ دسی‌بل رسید. بهبود نسبت به دو‌مکعبی بین ۱٫۵۱۳ تا ۲٫۲۸۳ دسی‌بل بود؛ بنابراین شبکه فقط روی مجموعهٔ کوچک \lr{Set5} بهتر نشده و روی تصاویر متنوع‌تر نیز مزیت پایدار دارد.

در \lr{B100} کمترین افزایش PSNR نسبت به دو‌مکعبی، یعنی ۱٫۹۴۲، ۱٫۲۵۱ و ۱٫۰۰۷ دسی‌بل، دیده شد. این مجموعه بیشتر شامل تصاویر طبیعی است و بافت‌های نامنظم آن نسبت به خطوط تکراری شهری و کمیک، سود کمتری از بازسازی ساختاری شبکه می‌برند. در عوض فاصلهٔ ما با مقاله در همین مجموعه از همه کمتر بود: ۰٫۶۱۵، ۰٫۶۵۰ و ۰٫۶۲۲ دسی‌بل.

در \lr{Urban100} شروع گرم نقش مهم‌تری داشت و برای مقیاس‌های ۲، ۳ و ۴ به‌ترتیب ۰٫۶۷۵، ۰٫۷۱۶ و ۰٫۴۵۸ دسی‌بل افزود. خروجی نهایی ۳۰٫۳۴۳، ۲۶٫۶۵۵ و ۲۴٫۷۸۱ دسی‌بل بود. خطوط نما، پنجره‌ها و الگوهای تکراری ساختمان‌ها با میدان دید ۱۱ در ۱۱ پیمانه‌سازی و توجه فضایی معماری سازگارند، ولی فاصلهٔ باقی‌مانده با مقاله نشان می‌دهد یادگیری همین جزئیات به آموزش طولانی‌تر حساس است.

در \lr{Manga109} هم PSNR و هم SSIM بیشترین سود را از مدل گرفتند. خروجی نهایی به ۳۶٫۸۶۷/۰٫۹۷۲۰، ۳۱٫۲۶۹/۰٫۹۲۳۶ و ۲۸٫۱۶۸/۰٫۸۶۹۷ رسید. افزایش شروع گرم در مقیاس‌های سه و چهار به‌ترتیب ۱٫۴۰۸ و ۱٫۰۹۲ دسی‌بل بود؛ این دو مقدار بزرگ‌ترین اثر شروع گرم در کل جدول‌اند. در نتیجه، آموزش دوم بیش از همه برای بازسازی خطوط باریک و بافت‌های منظم کمیک مؤثر بوده است.

جدول \ref{tab:our-psnr-differences} دو مقایسهٔ مهم را جدا می‌کند. ستون نخست نشان می‌دهد شروع گرم در همهٔ حالت‌ها مفید بوده است. ستون دوم فاصلهٔ خروجی نهایی با خط پایه را نشان می‌دهد. ستون سوم اختلاف با مقاله است؛ منفی بودن آن یعنی اجرای ۲۰۰‌دوره‌ای ما هنوز به آموزش طولانی مقاله نرسیده است.

\begin{table}[H]
\centering
\caption{اختلاف PSNR خروجی نهایی برحسب دسی‌بل؛ «مرحلهٔ اول» خروجی ۱۰۰‌دوره‌ای پیش از شروع گرم است}
\label{tab:our-psnr-differences}
\resizebox{0.88\textwidth}{!}{%
\begin{tabular}{clrrr}
\toprule
مقیاس & مجموعه & نسبت به مرحلهٔ اول & نسبت به دو‌مکعبی & نسبت به مقاله \\
\midrule
۲ & \lr{Set5}     & +۰٫۴۰۴ & +۳٫۲۰۳ & -۱٫۰۲۷ \\
۲ & \lr{Set14}    & +۰٫۳۴۵ & +۲٫۲۸۳ & -۰٫۹۵۴ \\
۲ & \lr{B100}     & +۰٫۲۹۸ & +۱٫۹۴۲ & -۰٫۶۱۵ \\
۲ & \lr{Urban100} & +۰٫۶۷۵ & +۳٫۲۷۱ & -۲٫۳۳۷ \\
۲ & \lr{Manga109} & +۰٫۷۹۲ & +۵٫۶۱۵ & -۲٫۲۶۳ \\
۳ & \lr{Set5}     & +۰٫۷۱۳ & +۲٫۵۶۶ & -۱٫۴۴۱ \\
۳ & \lr{Set14}    & +۰٫۳۵۱ & +۱٫۶۷۹ & -۱٫۰۷۴ \\
۳ & \lr{B100}     & +۰٫۳۱۰ & +۱٫۲۵۱ & -۰٫۶۵۰ \\
۳ & \lr{Urban100} & +۰٫۷۱۶ & +۲٫۰۷۳ & -۱٫۹۳۵ \\
۳ & \lr{Manga109} & +۱٫۴۰۸ & +۴٫۰۷۱ & -۲٫۸۴۱ \\
۴ & \lr{Set5}     & +۰٫۵۲۵ & +۲٫۱۷۷ & -۱٫۶۸۱ \\
۴ & \lr{Set14}    & +۰٫۳۷۳ & +۱٫۵۱۳ & -۱٫۰۳۷ \\
۴ & \lr{B100}     & +۰٫۲۶۵ & +۱٫۰۰۷ & -۰٫۶۲۲ \\
۴ & \lr{Urban100} & +۰٫۴۵۸ & +۱٫۵۴۰ & -۱٫۶۰۹ \\
۴ & \lr{Manga109} & +۱٫۰۹۲ & +۳٫۰۹۶ & -۲٫۷۶۲ \\
\bottomrule
\end{tabular}}
\end{table}

شکل \ref{fig:warm-start-gain} افزایش PSNR مرحلهٔ شروع گرم نسبت به مرحلهٔ نخست را برای همهٔ ۱۵ حالت نشان می‌دهد. همهٔ میله‌ها مثبت‌اند؛ بیشترین افزایش مربوط به \lr{Manga109} در مقیاس سه است.

\begin{figure}[H]
    \centering
    \includegraphics[width=0.9\textwidth]{images/warm_start_psnr_gain.png}
    \caption{افزایش PSNR بر اثر مرحلهٔ ۱۰۰‌دوره‌ای شروع گرم نسبت به پایان مرحلهٔ نخست، برای هر مجموعه و مقیاس.}
    \label{fig:warm-start-gain}
\end{figure}

فاصله با مقاله روی \lr{B100} کمتر و روی مجموعه‌های ساختاری \lr{Urban100} و \lr{Manga109} بیشتر است. برداشت ما این است که بافت‌ها و وابستگی‌های پیچیده‌تر برای استفادهٔ کامل از میدان دید بزرگ شبکه به تعداد به‌روزرسانی بیشتری نیاز دارند. این یک استنباط از الگوی جدول است، نه آزمایش حذف مستقل. همچنین در مقیاس چهار، بهترین نقطهٔ وارسی پایشی \lr{Set5} به ۳۰٫۹۶۷ دسی‌بل رسید، اما برای رعایت قاعدهٔ یکسان و پرهیز از انتخاب روی آزمون، مقدار پایان دوره یعنی ۳۰٫۷۷۹ گزارش شد.

\section{زمان استنتاج و بررسی کیفی}

جدول \ref{tab:our-inference-time} میانگین زمان اجرای یک تصویر را روی همان کارت T4 نشان می‌دهد. زمان به اندازه و مقیاس تصویر وابسته است؛ ازاین‌رو مقایسهٔ مستقیم اعداد مجموعه‌های متفاوت یا مقایسه با سخت‌افزار مقاله منصفانه نیست. الگوی کلی درست است: با بزرگ‌تر شدن ورودی کم‌وضوح، زمان افزایش می‌یابد و \lr{Manga109} و \lr{Urban100} سنگین‌ترین مجموعه‌ها هستند.

\begin{table}[H]
\centering
\caption{میانگین زمان استنتاج هر تصویر روی کارت \lr{Tesla T4}، برحسب میلی‌ثانیه}
\label{tab:our-inference-time}
\begin{tabular}{rrrrrr}
\toprule
مقیاس & \lr{Set5} & \lr{Set14} & \lr{B100} & \lr{Urban100} & \lr{Manga109} \\
\midrule
۲ & ۶۰٫۵ & ۱۱۴٫۶ & ۷۳٫۰ & ۴۰۱٫۴ & ۵۱۶٫۷ \\
۳ & ۳۰٫۸ & ۶۰٫۶ & ۳۹٫۸ & ۱۸۲٫۲ & ۲۲۷٫۵ \\
۴ & ۲۰٫۶ & ۴۷٫۰ & ۲۸٫۵ & ۹۸٫۰ & ۱۲۸٫۴ \\
\bottomrule
\end{tabular}
\end{table}

\subsection{نمونه‌های کیفی در حالت‌های مختلف}

برای اینکه بررسی بصری به یک تصویر شهری محدود نماند، هفت نمونه از پنج مجموعه و هر سه ضریب بزرگ‌نمایی در شکل‌های \ref{fig:qual-set5-baby-x2} تا \ref{fig:qual-urban-x4} آورده‌ایم. هر چهار پنل مستقیماً از آرایه‌های ورودی، درون‌یابی، خروجی شبکه و مرجع ساخته شده‌اند و تصویر صفحه نیستند. این نمونه‌ها برای پوشش‌دادن چهره و بافت پارچه، بافت طبیعی، خطوط معماری، تصویر کمیک و بزرگ‌نمایی دشوار چهاربرابری انتخاب شده‌اند. همهٔ تصاویر مقایسه‌ای تولیدشده، از جمله نمونه‌های دیگری که در متن جا نشده‌اند، در پوشهٔ \lr{benchmarks} هر اجرا موجودند.

\begin{table}[H]
\centering
\caption{سنجه‌های عددی نمونه‌های کیفی نمایش‌داده‌شده؛ اعداد از ردیف همان تصویر در فایل‌های CSV گرفته شده‌اند}
\label{tab:qualitative-samples}
\resizebox{0.95\textwidth}{!}{%
\begin{tabular}{cllrrrr}
\toprule
مقیاس & مجموعه & تصویر & PSNR مدل & SSIM مدل & PSNR دو‌مکعبی & SSIM دو‌مکعبی \\
\midrule
۲ & \lr{Set5} & \lr{baby} & ۳۸٫۶۷۵ & ۰٫۹۶۶۴ & ۳۷٫۲۳۶ & ۰٫۹۵۴۹ \\
۲ & \lr{Set14} & \lr{barbara} & ۲۸٫۳۴۲ & ۰٫۸۷۶۵ & ۲۸٫۱۰۷ & ۰٫۸۴۶۷ \\
۲ & \lr{Urban100} & \lr{img\_001} & ۳۱٫۷۲۹ & ۰٫۸۹۴۹ & ۲۸٫۳۰۶ & ۰٫۸۲۶۱ \\
۳ & \lr{B100} & \lr{101085} & ۲۴٫۷۹۳ & ۰٫۶۵۲۳ & ۲۴٫۰۴۱ & ۰٫۵۹۰۷ \\
۳ & \lr{Manga109} & \lr{ARMS} & ۲۹٫۵۱۸ & ۰٫۸۵۲۸ & ۲۷٫۰۷۶ & ۰٫۷۷۵۴ \\
۴ & \lr{Set5} & \lr{bird} & ۳۳٫۰۷۹ & ۰٫۹۲۱۶ & ۳۰٫۴۴۰ & ۰٫۸۷۷۶ \\
۴ & \lr{Urban100} & \lr{img\_002} & ۲۵٫۶۱۰ & ۰٫۷۷۰۹ & ۲۴٫۲۷۷ & ۰٫۶۸۲۸ \\
\bottomrule
\end{tabular}}
\end{table}

در نمونهٔ \lr{baby}، مدل بافت کلاه، مرز چشم‌ها و خطوط ظریف صورت را بدون ایجاد لبهٔ شدید حفظ کرده و ۱٫۴۳۹ دسی‌بل از دو‌مکعبی بهتر است. این تصویر نمونه‌ای از ناحیه‌های نرم همراه با چند بافت ریز است.

\begin{figure}[H]
    \centering
    \includegraphics[width=\textwidth]{images/qualitative_set5_baby_x2.png}
    \caption{نمونهٔ \lr{baby} از \lr{Set5} در مقیاس دو پس از شروع گرم؛ از راست به چپ: مرجع پروضوح، STSN، دو‌مکعبی و ورودی کم‌وضوح.}
    \label{fig:qual-set5-baby-x2}
\end{figure}

نمونهٔ \lr{barbara} به‌خاطر پارچه‌های راه‌راه و بافت‌های بسامدبالا دشوار است. افزایش PSNR فقط ۰٫۲۳۴ دسی‌بل است، اما SSIM از ۰٫۸۴۶۷ به ۰٫۸۷۶۵ رسیده؛ یعنی بهبود ساختار محلی از کاهش خطای پیکسلی مشهودتر بوده است.

\begin{figure}[H]
    \centering
    \includegraphics[width=\textwidth]{images/qualitative_set14_barbara_x2.png}
    \caption{نمونهٔ \lr{barbara} از \lr{Set14} در مقیاس دو؛ بافت پارچه و خطوط ریز میز، تفاوت بازسازی ساختاری را آشکار می‌کنند. ترتیب پنل‌ها از راست: مرجع، STSN، دو‌مکعبی و ورودی است.}
    \label{fig:qual-set14-barbara-x2}
\end{figure}

در شکل \ref{fig:our-qualitative}، تکرار پنجره‌ها و لبه‌های ساختمان نسبت به دو‌مکعبی منظم‌تر شده و افزایش ۳٫۴۲۳ دسی‌بلی با این تفاوت بصری سازگار است. بااین‌حال، در جزئیات بسیار ریز هنوز اختلاف با مرجع دیده می‌شود.

\begin{figure}[H]
    \centering
    \includegraphics[width=\textwidth]{images/comparison_x2_urban100_warm.png}
    \caption{مقایسهٔ کیفی یک تصویر \lr{Urban100} در مقیاس دو پس از شروع گرم؛ از راست به چپ: مرجع پروضوح، خروجی STSN، درون‌یابی دو‌مکعبی و ورودی کم‌وضوح. شکل خروجی مستقیم برنامه است.}
    \label{fig:our-qualitative}
\end{figure}

نمونهٔ \lr{101085} از \lr{B100} بافت طبیعی و نامنظم دارد. در این‌جا افزایش PSNR برابر ۰٫۷۵۲ دسی‌بل است، ولی SSIM حدود ۰٫۰۶۲ افزایش یافته است. مدل مرز پیکره‌ها را واضح‌تر می‌کند، اما نبود الگوی تکراری منظم باعث می‌شود سود آن از نمونه‌های شهری و کمیک کمتر باشد.

\begin{figure}[H]
    \centering
    \includegraphics[width=\textwidth]{images/qualitative_b100_101085_x3.png}
    \caption{نمونهٔ \lr{101085} از \lr{B100} در مقیاس سه؛ از راست به چپ: مرجع، STSN، دو‌مکعبی و ورودی کم‌وضوح. این نمونه رفتار مدل روی بافت طبیعی نامنظم را نشان می‌دهد.}
    \label{fig:qual-b100-x3}
\end{figure}

در تصویر \lr{ARMS} خطوط باریک نقشه، نوشته‌ها و مرزهای ترسیمی فراوان‌اند. STSN نسبت به دو‌مکعبی ۲٫۴۴۲ دسی‌بل و حدود ۰٫۰۷۷ SSIM بهتر است. این نتیجه با بهبود بالای میانگین \lr{Manga109} هم‌جهت است و نشان می‌دهد مدل در بازسازی ساختارهای خطی و تکرارشونده موفق‌تر است.

\begin{figure}[H]
    \centering
    \includegraphics[width=\textwidth]{images/qualitative_manga109_arms_x3.png}
    \caption{نمونهٔ \lr{ARMS} از \lr{Manga109} در مقیاس سه؛ خروجی STSN خطوط ترسیمی و نوشته‌های ریز را نسبت به دو‌مکعبی واضح‌تر نگه می‌دارد.}
    \label{fig:qual-manga-x3}
\end{figure}

برای نمایش دشواری مقیاس چهار، دو محتوای متفاوت انتخاب شد. در تصویر \lr{bird}، مرز منقار و چشم و شاخه‌های باریک واضح‌تر شده‌اند و افزایش PSNR به ۲٫۶۳۹ دسی‌بل می‌رسد؛ بااین‌حال، فاصله با مرجع هنوز بیشتر از نمونه‌های مقیاس دو است.

\begin{figure}[H]
    \centering
    \includegraphics[width=\textwidth]{images/qualitative_set5_bird_x4.png}
    \caption{نمونهٔ \lr{bird} از \lr{Set5} در بزرگ‌نمایی چهاربرابری؛ این حالت اثر کاهش شدید اطلاعات ورودی بر مرزهای طبیعی را نشان می‌دهد. ترتیب پنل‌ها از راست: مرجع، STSN، دو‌مکعبی و ورودی است.}
    \label{fig:qual-set5-bird-x4}
\end{figure}

نمونهٔ شهری مقیاس چهار شامل خط‌های مورب و منحنی‌های تکرارشوندهٔ سازه است. مدل ۱٫۳۳۳ دسی‌بل و حدود ۰٫۰۸۸ SSIM بهتر از دو‌مکعبی عمل می‌کند و پیوستگی تیرها را بهتر نگه می‌دارد، ولی در بخش‌های بسیار متراکم بخشی از خطوط ظریف همچنان نرم می‌شوند.

\begin{figure}[H]
    \centering
    \includegraphics[width=\textwidth]{images/qualitative_urban100_img002_x4.png}
    \caption{نمونهٔ \lr{img\_002} از \lr{Urban100} در مقیاس چهار؛ خطوط مورب و قوس‌های تکراری سازه برای مقایسهٔ بازسازی هندسی مناسب‌اند.}
    \label{fig:qual-urban-x4}
\end{figure}

\iffalse
\section{تنظیمات اجرای سریع ما}

آموزش کامل مقاله روی سخت‌افزار در دسترس و زمان تحویل پروژه عملی نبود. برای آزمون صحت، هشت تصویر $128\times128$ تولید و چهار تصویر مستقل با همان مولد برای ارزیابی ساخته شد. شبکه همچنان مدل کامل ۵۰ کاناله با چهار گروه و چهار بلوک بود. اندازهٔ وصلهٔ HR برابر ۴۸، اندازهٔ دسته ۲، تعداد دوره ۲۰ و تعداد دسته در هر دوره ۸ بود؛ در مجموع فقط ۱۶۰ به‌روزرسانی انجام شد. اجرا روی کارت en{NVIDIA GeForce GTX 1050} و en{PyTorch 2.7.1} انجام شد.

جدول \ref{tab:quick-training} روند همگرایی را در چند دورهٔ منتخب نشان می‌دهد. زیان آموزش از $0.8057$ به حدود $0.06$ کاهش یافت و PSNR اعتبارسنجی تا دورهٔ ۱۹ به $26.79$ دسی‌بل رسید. نوسان دوره‌های پایانی ناشی از اندازهٔ بسیار کوچک داده و نرخ یادگیری ثابت در این بازه است.

\begin{table}[H]
\centering
\caption{روند اجرای سریع مدل کامل روی دادهٔ کوچک؛ اعداد مستقیماً از فایل تاریخچهٔ اجرای کد گرفته شده‌اند}
\label{tab:quick-training}
\begin{tabular}{rrrr}
\toprule
دوره & زیان $L_1$ آموزش & PSNR اعتبارسنجی & SSIM اعتبارسنجی \\
\midrule
۱ & ۰٫۸۰۵۷ & ۱۱٫۸۹۰ & ۰٫۰۲۵۹ \\
۴ & ۰٫۱۸۰۰ & ۱۹٫۰۰۳ & ۰٫۳۰۷۶ \\
۸ & ۰٫۰۹۴۴ & ۲۲٫۸۹۳ & ۰٫۵۲۶۳ \\
۱۲ & ۰٫۰۷۶۲ & ۲۵٫۰۶۳ & ۰٫۶۱۴۰ \\
۱۶ & ۰٫۰۶۳۳ & ۲۵٫۸۷۵ & ۰٫۶۱۲۱ \\
۱۹ & ۰٫۰۵۷۲ & ۲۶٫۷۸۷ & ۰٫۶۸۶۶ \\
۲۰ & ۰٫۰۶۱۴ & ۲۶٫۳۵۵ & ۰٫۶۳۲۹ \\
\bottomrule
\end{tabular}
\end{table}

شکل \ref{fig:training-curve} روند کامل را نشان می‌دهد. نقطهٔ وارسی دورهٔ ۱۹ به‌عنوان بهترین مدل ذخیره شد؛ بنابراین سنجهٔ نهایی از همان دوره و نه آخرین دوره گزارش شده است.

\begin{figure}[H]
    \centering
    \includegraphics[width=0.82\textwidth]{images/training_curve.png}
    \caption{تغییر زیان قدرمطلق آموزش و PSNR اعتبارسنجی در اجرای سریع ۲۰‌دوره‌ای؛ هر دوره شامل هشت دسته است.}
    \label{fig:training-curve}
\end{figure}

\section{نتایج عددی اجرای سریع}

نتیجهٔ هر چهار تصویر در جدول \ref{tab:quick-images} آمده است. میانگین زمان استنتاج $53.38$ میلی‌ثانیه بود. این زمان برای ورودی کوچک $64\times64$ اندازه‌گیری شده و نباید با زمان ۶۴۰ میلی‌ثانیهٔ مقاله روی تصاویر اعتبارسنجی en{DIV2K} مقایسهٔ مستقیم شود.

\begin{table}[H]
\centering
\caption{نتایج هر تصویر در اجرای سریع؛ ارزیابی روی کانال Y و پس از حذف حاشیهٔ دو پیکسلی انجام شده است}
\label{tab:quick-images}
\begin{tabular}{lrrrr}
\toprule
تصویر & PSNR مدل & SSIM مدل & PSNR دو‌مکعبی & SSIM دو‌مکعبی \\
\midrule
اعتبارسنجی ۱ & ۲۶٫۸۳۷ & ۰٫۶۶۳۱ & ۳۳٫۶۸۰ & ۰٫۹۴۹۲ \\
اعتبارسنجی ۲ & ۲۶٫۶۲۵ & ۰٫۶۹۸۷ & ۳۳٫۷۰۵ & ۰٫۹۵۵۷ \\
اعتبارسنجی ۳ & ۲۶٫۱۸۹ & ۰٫۶۹۴۲ & ۳۱٫۷۲۳ & ۰٫۹۴۵۴ \\
اعتبارسنجی ۴ & ۲۷٫۴۹۹ & ۰٫۶۹۰۶ & ۳۴٫۹۱۶ & ۰٫۹۵۵۵ \\
\midrule
میانگین & ۲۶٫۷۸۷ & ۰٫۶۸۶۶ & ۳۳٫۵۰۶ & ۰٫۹۵۱۵ \\
\bottomrule
\end{tabular}
\end{table}

دو‌مکعبی در این آزمایش $6.72$ دسی‌بل بهتر از مدل است. این نتیجه شکست معماری مقاله نیست؛ مدل مقاله ۱۰۰۰ دوره روی ۸۰۰ تصویر طبیعی و وصله‌های بسیار بیشتر آموزش دیده است، در حالی که این اجرا تنها ۱۶۰ گام روی تصویرهای ساختگی دارد. خط پایهٔ دو‌مکعبی از قبل پاسخ هموار مناسبی برای الگوهای سینوسی دادهٔ ما دارد، اما لایهٔ بازآرایی زیرپیکسلی شبکه باید نگاشت خود را از وزن تصادفی بیاموزد. در شکل \ref{fig:quick-comparison} آثار شبکه‌ای خروجی STSN هنوز دیده می‌شود و همین موضوع با SSIM پایین‌تر سازگار است. درج این شکل برای نشان دادن محدودیت اجرای کوتاه است، نه ادعای برتری بصری.

\begin{figure}[H]
    \centering
    \includegraphics[width=\textwidth]{images/quick_comparison.png}
    \caption{مقایسهٔ کیفی نمونهٔ اول اجرای سریع، از راست به چپ: مرجع پروضوح، خروجی STSN، درون‌یابی دو‌مکعبی و ورودی کم‌وضوح. تصویرها خروجی مستقیم برنامه‌اند و تصویر صفحه نیستند.}
    \label{fig:quick-comparison}
\end{figure}

\fi
\section{نسخهٔ مسابقهٔ NTIRE 2023}

مقاله در بخش پایانی نسخهٔ بزرگ‌تری را برای مسابقهٔ \lr{NTIRE 2023} توضیح می‌دهد. در آن نسخه تعداد گروه‌ها ۵، تعداد بلوک هر گروه ۴، تعداد کانال ویژگی ۱۵۰ و کانال میانی توجه فضایی ۳۲ است. بازهٔ RGB نیز به‌جای یک، ۲۵۵ در نظر گرفته شده است. آموزش اولیه با \lr{DIV2K} و \lr{LSDIR}، وصلهٔ ۱۹۲، دستهٔ ۱۶ و ۷۰ دوره انجام شده؛ سپس شروع گرم ۴۵۰‌دوره‌ای و در مرحلهٔ آخر آموزش ۲۰۰‌دوره‌ای با \lr{DIV2K} و \lr{Flickr2K} و نرخ یادگیری $5\times10^{-5}$ اجرا شده است. این تیم با PSNR برابر ۳۰٫۷۸۰ و SSIM برابر ۰٫۸۵۸۲ رتبهٔ نهم جدول مقاله را به‌دست آورده است. کلاس مدل ما همین تغییر تعداد گروه، کانال و کانال توجه را می‌پذیرد، ولی وزن نسخهٔ مسابقه به‌علت نبود مجموعه‌های آموزشی افزوده در محدودهٔ این پروژه آموزش داده نشده است.

\chapter{تحلیل، محدودیت‌ها و بازتولیدپذیری}

\section{چرا تطبیق پارامتر برای ما مهم بود؟}

فقط برابر بودن شکل خروجی برای اثبات بازپیاده‌سازی معماری کافی نیست. برای مثال، تغییر کرنل عمقی از ۱۱ به ۷ همچنان خروجی با ابعاد درست می‌دهد ولی رفتار و ظرفیت مدل متفاوت می‌شود. تطبیق هفت عدد مستقل در جدول \ref{tab:param-ablation} قید قوی‌تری ایجاد می‌کند: سهم هر پیمانه‌سازی، پرسپترون، توجه فضایی، کانولوشن گروه و فرافکنی به‌طور جداگانه سنجیده شده است. به همین دلیل اطمینان داریم پیاده‌سازی از نظر ساختار پارامتری با مدل مقاله سازگار است.

\section{محدودیت‌های مقایسه}

سه محدودیت باید روشن باشد. نخست، مقاله وزن‌های آموزش‌دیده را منتشر نکرده و برخی جزئیات معماری در متن صریح نیستند؛ ما آن‌ها را از شمار پارامترها استنتاج و فرض‌ها را مستند کرده‌ایم. دوم، روش دقیق کاهش دو‌مکعبی در کتابخانه‌های مختلف اندکی تفاوت دارد. برای آموزش از جفت‌های رسمی DIV2K و برای ارزیابی از جفت‌های آماده و بازنمونه‌برداری دو‌مکعبی متلب استفاده کردیم تا این اختلاف کم شود. سوم، کامل بودن داده به معنای برابر بودن بودجهٔ بهینه‌سازی نیست: مقاله ۱۰۰۰ دوره و سپس بازآموزی را گزارش می‌کند، در حالی که مطابق محدودیت پروژه هر مرحلهٔ ما ۱۰۰ دوره بود. بنابراین اختلاف ستون آخر جدول \ref{tab:our-psnr-differences} را به‌عنوان نتیجهٔ بازپیاده‌سازی ۲۰۰‌دوره‌ای گزارش می‌کنیم، نه بازتولید کامل بودجهٔ محاسباتی مقاله.

\section{روش اجرای پروژه}

پس از نصب وابستگی‌های فایل \lr{requirements.txt}، دستور زیر مسیر کوچک و سریع را برای کنترل سلامت پیاده‌سازی انجام می‌دهد:

\begin{latin}
\begin{verbatim}
python scripts/main.py --quick
\end{verbatim}
\end{latin}

برای بررسی عددی معماری می‌توان نوشت:

\begin{latin}
\begin{verbatim}
python scripts/ablation.py
python -m unittest discover -s tests -v
\end{verbatim}
\end{latin}

برای تکرار آزمایش گزارش‌شده روی هر سه مقیاس، پس از قرار دادن جفت‌های DIV2K و مجموعه‌های معیار در مسیرهای نسبی پیکربندی، دستور زیر هر دو مرحلهٔ ۱۰۰‌دوره‌ای را اجرا می‌کند:

\begin{latin}
\begin{verbatim}
python scripts/run_full_experiment.py --epochs 100
\end{verbatim}
\end{latin}

خروجی‌ها شامل بهترین و آخرین نقطهٔ وارسی هر مرحله، تاریخچهٔ JSON، جدول CSV هر تصویر، نمودار آموزش، مقایسهٔ کیفی و جدول تجمیعی پنج مجموعه هستند. بذر تصادفی پیش‌فرض ۸۱۰۱۰۴۰۳۹ است. آزمون‌های خودکار سه مورد را بررسی کردند: برابری پارامترهای هفت گونهٔ مدل، شکل خروجی همهٔ ضرایب مقیاس و هم‌ترازی وصله‌های LR/HR؛ هر سه آزمون موفق بودند. اسکریپت \lr{collect\_kaggle\_results.py} نیز کامل بودن شش اجرا را کنترل و جدول‌های گزارش را از فایل‌های خام تولید می‌کند.

\section{منابع کد و دادهٔ بیرونی}

بخش اصلی مدل، حلقهٔ آموزش، سنجه‌ها و شکل‌ها در این پروژه نوشته شده‌اند و کد آماده‌ای از اینترنت در آن‌ها کپی نشده است. برای فهم عملگر مبنا، مقاله و مخزن رسمی Conv2Former بررسی شدند \cite{Hou2022,Conv2FormerGitHub}. تصاویر DIV2K از نسخه‌های آینه‌ای دادهٔ رسمی در کگل و مجموعه‌های آزمون از منابع معرفی‌شده در کد دریافت شدند؛ فایل‌های کم‌وضوح سه مجموعه نیز با اسکریپت متلب همراه پروژه ساخته شدند. نشانی دریافت Set5 داخل \lr{download\_set5.py} و در منبع \cite{Set5HF} آمده است. همهٔ تصاویر گزارش خروجی مستقیم کد یا نمودار ترسیم‌شده از تاریخچهٔ عددی‌اند و هیچ تصویر صفحه‌ای در گزارش استفاده نشده است.

\chapter{نتیجه‌گیری}

در این پروژه کل روش STSN را از مرحلهٔ داده تا گزارش نتیجه بازپیاده‌سازی کردیم. ایدهٔ مهم روش این است که می‌توان ساختار مانده و پیش‌نرمال‌سازی ترنسفورمر را نگه داشت، ولی خودتوجهی درجه‌دوم را با پیمانه‌سازی کانولوشنی خطی جایگزین کرد. کانولوشن $3\times3$ پایانی ویژگی محلی، کرنل عمقی ۱۱ میدان دید، پرسپترون اختلاط کانال و توجه فضایی انتخاب مکان‌های مهم را بر عهده دارند.

قوی‌ترین راستی‌آزمایی ساختاری ما برابری دقیق شمار پارامتر مدل کامل و همهٔ گونه‌های حذف مؤلفه با مقاله بود. راستی‌آزمایی تجربی نیز روی کل ۸۰۰ تصویر DIV2K، سه مقیاس و همهٔ تصاویر پنج مجموعهٔ معیار انجام شد. مدل نهایی در هر ۱۵ حالت از دو‌مکعبی بهتر بود و شروع گرم نیز در هر ۱۵ حالت نسبت به مرحلهٔ نخست PSNR را افزایش داد. بهترین بهبود نسبت به دو‌مکعبی ۵٫۶۱۵ دسی‌بل روی Manga109 در مقیاس دو بود. با وجود این، نتیجه هنوز در همهٔ حالت‌ها پایین‌تر از مقاله است؛ اختلاف ۰٫۶۱۵ تا ۲٫۸۴۱ دسی‌بل با کاهش بودجه از آموزش ۱۰۰۰‌دوره‌ای و بازآموزی مقاله به دو مرحلهٔ ۱۰۰‌دوره‌ای ما سازگار است. کد و دادهٔ ذخیره‌شده امکان ادامه تا بودجهٔ مقاله را بدون تغییر معماری فراهم می‌کنند.

\renewcommand{\bibname}{منابع}
\begin{thebibliography}{99}
\begin{latin}
\bibitem{Gendy2023}
G. Gendy, N. Sabor, J. Hou, and G. He, ``A Simple Transformer-Style Network for Lightweight Image Super-Resolution,'' in \textit{CVPR Workshops}, pp. 1484--1494, 2023.

\bibitem{Hou2022}
Q. Hou, C.-Z. Lu, M.-M. Cheng, and J. Feng, ``Conv2Former: A Simple Transformer-Style ConvNet for Visual Recognition,'' arXiv:2211.11943, 2022.

\bibitem{Liu2020}
J. Liu, W. Zhang, Y. Tang, J. Tang, and G. Wu, ``Residual Feature Aggregation Network for Image Super-Resolution,'' in \textit{CVPR}, pp. 2359--2368, 2020.

\bibitem{Agustsson2017}
E. Agustsson and R. Timofte, ``NTIRE 2017 Challenge on Single Image Super-Resolution: Dataset and Study,'' in \textit{CVPR Workshops}, pp. 126--135, 2017.

\bibitem{Bevilacqua2012}
M. Bevilacqua, A. Roumy, C. Guillemot, and M.-L. Alberi-Morel, ``Low-Complexity Single-Image Super-Resolution Based on Nonnegative Neighbor Embedding,'' in \textit{BMVC}, 2012.

\bibitem{Wang2004}
Z. Wang, A. C. Bovik, H. R. Sheikh, and E. P. Simoncelli, ``Image Quality Assessment: From Error Visibility to Structural Similarity,'' \textit{IEEE Transactions on Image Processing}, vol. 13, no. 4, pp. 600--612, 2004.

\bibitem{Kingma2014}
D. P. Kingma and J. Ba, ``Adam: A Method for Stochastic Optimization,'' arXiv:1412.6980, 2014.

\bibitem{Shi2016}
W. Shi et al., ``Real-Time Single Image and Video Super-Resolution Using an Efficient Sub-Pixel Convolutional Neural Network,'' in \textit{CVPR}, pp. 1874--1883, 2016.

\bibitem{Conv2FormerGitHub}
HVision-NKU, ``Official PyTorch Implementation of Conv2Former,'' \url{https://github.com/HVision-NKU/Conv2Former}, accessed August 2026.

\bibitem{Set5HF}
E. Siow, ``Set5 Dataset Card and Image Archives,'' \url{https://huggingface.co/datasets/eugenesiow/Set5}, accessed August 2026.
\end{latin}
\end{thebibliography}
